This paper has demonstrated that UX design decisions have measurable environmental consequences 
that scale across billions of users. The ethical analysis established that current practices violate principles of 
climate justice and autonomy. Distributive injustice occurs when users' 
convenience is subsidized by climate impacts on vulnerable populations. Autonomy violations occur when users are denied 
information necessary to make choices aligned with their own values. The "black 
box" problem, where companies do not disclose energy consumption, creates 
an injustice that prevents both users and designers from making informed 
decisions. Designers operate within structural constraints that make individual 
ethical action insufficient. Performance evaluations based on engagement metrics 
rather than environmental impact create professional incentives that systematically 
punish sustainable design choices. While designers retain agency to advocate for 
eco-friendly alternatives, they cannot 
overcome business models that require growth maximization.

\par The normative conclusion is therefore twofold. First, designers bear professional 
responsibility to minimize environmental impact within their sphere of influence: 
being aware, learning about and advocating for sustainable alternatives, and supporting 
regulatory efforts. Second, systemic solutions are necessary to make ethical 
design professionally viable: mandatory energy disclosure enables informed choice, 
while regulatory default standards shift competitive dynamics to reward 
sustainability rather than punish it.

\par As digital services continue expanding and AI adoption accelerates, the question 
of who bears responsibility for environmental consequences will intensify. The 
framework developed here, combining individual professional responsibility with 
recognition of structural barriers requiring regulatory intervention, provides a 
path forward that avoids both naive individualism and fatalistic determinism. 
Designers must act within their sphere while advocating for systemic change that 
expands what ethical action can accomplish.